\section{Conclusiones}
\label{sec:conclusiones}

En este trabajo desarrollamos y evaluamos un prototipo de clasificación de enfermedades
foliares en pimiento, papa y tomate a partir de modelos livianos preentrenados en ImageNet,
entrenados íntegramente en CPU. Las principales conclusiones son las siguientes:

\begin{enumerate}
  \item \textbf{El fine-tuning parcial es más determinante que la elección del backbone.} Las
        representaciones de ImageNet ya alcanzan un macro-F1 de 0,96 con un linear probe,
        con diferencias mínimas entre arquitecturas móviles. Reentrenar solo las últimas
        etapas lleva el macro-F1 a 0,987--0,993, con un costo de entrenamiento compatible con
        una CPU de notebook. La comparación también cambia la resolución y el clasificador.
  \item \textbf{Las arquitecturas móviles ofrecen el mejor compromiso.} Superan a ResNet18
        con muchos menos parámetros y permiten inferencia en tiempo real sin GPU.
        EfficientNet-B0 obtiene el mejor desempeño (17 errores en 2\,949 imágenes de test,
        17,5\,ms por imagen), mientras que MobileNetV3-Small es 2,5 veces más rápido a cambio
        de casi el doble de errores, lo que lo convierte en un candidato para dispositivos más limitados.
  \item \textbf{El valor del data augmentation no se ve en el test de la misma distribución.}
        Con diferencias pequeñas de macro-F1 en el test original, el augmentation reduce la caída ante
        perturbaciones de hasta 0,23 a menos de 0,015. El desempeño en el test original por sí
        solo no permite anticipar la robustez ante estas perturbaciones.
  \item \textbf{La activación Grad-CAM se concentra principalmente en la hoja.} Según
        Grad-CAM, el 65,7\,\% de la activación recae sobre la hoja, con un cociente medio por
        imagen de 1,7 respecto de su área. La mayoría de los errores también presenta su
        máximo de activación dentro de la hoja, sin que esto establezca su causa. Cerca de un tercio de la activación
        cae fuera de la hoja, lo que deja abierta la posibilidad de que el modelo aproveche
        regularidades del fondo de PlantVillage.
  \item \textbf{Un desempeño cercano al 100\,\% en PlantVillage no garantiza utilidad en el
        campo.} Sobre imágenes reales de PlantDoc, el accuracy cae al 36,7\,\% y el modelo
        comete errores con alta confianza, incluso confundiendo la especie. El augmentation
        con transformaciones geométricas y de color no alcanza para cubrir esa brecha de
        dominio, por lo que el prototipo, en su estado actual, no debería usarse para
        diagnosticar en condiciones reales.
\end{enumerate}

Entre las limitaciones del estudio se encuentran el uso de una única semilla, la muestra
externa reducida (30 imágenes) y la fuga de datos que el agrupamiento por dHash no evita entre
fotografías rotadas de una misma hoja. La elección final orientada por perturbaciones
del test es exploratoria. Como trabajo futuro proponemos incorporar imágenes de
campo al entrenamiento o reemplazar el fondo de las hojas de PlantVillage, agregar un mecanismo
de detección de imágenes fuera de distribución que permita al sistema abstenerse ante casos
dudosos, repetir los experimentos con varias semillas y un conjunto externo más amplio, y
evaluar el modelo cuantizado directamente en un teléfono móvil.
